\section{The \instname\ data reduction flow.}

The overall data flow of the \instname\ pipeline is displayed in Figure \ref{fig:reduction_cascade}.

 
\begin{figure}[]
\centering
\includegraphics[width=0.9\textwidth]{../instrument_tutorial/figures/reduction_cascade.jpg}
\caption{The data reduction cascade of the \instname\ workflow.}
\label{fig:reduction_cascade}
\end{figure}

The reduction cascade is organized in tasks, which represent well-defined steps in the process. Tasks can be grouped
inside sub-workflows.
Each task runs a recipe; the detailed description of the algorithms,
input, outputs and recipe parameters used in each recipe are available
in the pipeline manual. Here, we present only the description of most
important features.

The EDPS workflow is designed to execute the tasks that deliver
the final reduced data for each dataset. It can be either the product of a single exposure, or the combination of
multiple exposures. Only calibrations needed by the selected scientific exposures are processed.

It is possible to set EDPS to perform the data reduction until a certain step of the reduction chain (e.g. to reduce
only standard stars, or only flat fields).
This is done by specifying the desired tasks in the field \textbf{Select reduction target} of the \textbf{Raw Data} tab.

The reduction steps of the workflow are listed below. Before starting the reduction,
the parameters of the recipes associated to each task can be configured by pressing the button 
\includegraphics{figures/configure_dataset.jpg} close to each dataset configuration.
See Section \ref{sec:configuration_overview} for more information.

The reduction steps are:

%%%%%%%%%%%%%%%%%%%%%%%%%%%%%%%%%

\subsection{Subworkflow: Dark Background}

Recipes:
\texttt{sph\_ird\_master\_dark}, 
\texttt{sph\_ird\_ins\_bg}, 
\texttt{sph\_ird\_sky\_bg}

This subworkflow generates the dark-current, instrumental-background, and
sky-background calibration products. It consists of the following tasks:

\begin{itemize}
    \item \textbf{irdis\_dark\_imaging}
    \item \textbf{irdis\_static\_bad\_pixel\_map}
    \item \textbf{irdis\_instrument\_background}
    \item \textbf{irdis\_sky\_background}
\end{itemize}

First, \textbf{irdis\_dark\_imaging} uses recipe \texttt{sph\_ird\_master\_dark} to create the master dark frame by
combining raw dark exposures using the selected collapse algorithm (typically
Clean Mean). 

The task \textbf{irdis\_static\_bad\_pixel\_map} runs the same recipe,
\texttt{sph\_ird\_master\_dark}, 

and exports the
\texttt{IRD\_STATIC\_BADPIXELMAP} product.

Next, \textbf{irdis\_instrument\_background} uses \texttt{sph\_ird\_ins\_bg} to generate the instrumental-background
calibration frame by combining raw background exposures using the selected
collapse algorithm. Unlike the master dark recipe, no bad-pixel map is created.

Finally, \textbf{irdis\_sky\_background} uses \texttt{sph\_ird\_sky\_bg} to produce the sky-background calibration frame
by combining raw sky-background exposures. As for the instrumental background
recipe, no bad-pixel map is generated.

The resulting \texttt{MASTER\_DARK}, \texttt{INS\_BG}, and
\texttt{SKY\_BG} products are intended for subtraction from raw science or calibration frames. They are neither flat-fielded nor corrected
using other dark or background calibrations.

The output products are multi-extension FITS files containing the calibration image, a bad-pixel map, a weight map recording the number of input frames contributing to each pixel, and an RMS map.

\paragraph{Customization}

Recipe parameters:

\begin{itemize}
\item \texttt{ird.sky\_bg.coll\_alg} / \texttt{ird.master\_dark.coll\_alg} / \texttt{ird.ins\_bg.coll\_alg}: Collapse algorithm used to combine the input frames
(\texttt{0} = Mean,
 \texttt{1} = Median,
 \texttt{2} = Clean Mean; default).
\end{itemize}

%%%%%%%%%%%%%%%%%%%%%%%%%%%%%%%%%%%%%%%%%%%%%%%%%%%%%%%%%%%%%%

\subsection{Master Flat}

Recipe: \texttt{sph\_ird\_instrument\_flat}

The task \textbf{irdis\_flat\_field} uses the pipeline recipe to generate the IRDIS detector flat field from narrow- or broadband lamp exposures acquired at different illumination
levels. 

Background inputs are optional. Raw dark frames with matching detector integration times (DITs) may be provided. 
If no suitable raw darks are available, EDPS may use the following calibration products in decreasing order of priority
\texttt{SKY\_BG}
$\rightarrow$
\texttt{INS\_BG}
$\rightarrow$
\texttt{MASTER\_DARK}.


\paragraph{Customization}

Recipe parameters:

\begin{itemize}
\item \texttt{ird.instrument\_flat.coll\_alg}:
Collapse algorithm
(\texttt{0} = Mean,
 \texttt{1} = Median,
 \texttt{2} = Clean Mean; default).
\item \texttt{ird.instrument\_flat.robust\_fits:}
Enables robust linear fitting. The default value is
\texttt{FALSE}. Setting this parameter to \texttt{TRUE} can improve rejection of outliers such as cosmic rays.

\end{itemize}

%%%%%%%%%%%%%%%%%%%%%%%%%%%%%%%%%%%%%%%%%%%%%%%%%%%%%%%%%%%%%%

\subsection{Distortion Map}

Recipe: \texttt{sph\_ird\_distortion\_map}

The task \textbf{irdis\_distortion\_map} uses the recipe to derive the IRDIS geometric distortion
map from calibration observations that have been reduced as standard
field-stabilized science images without dithering. The reduction may optionally
include dark subtraction and flat-field correction.

After combining the input frames, point sources are detected using a
user-defined detection threshold. If no reference point pattern is provided,
the recipe generates one automatically. Otherwise, the detected source
positions are compared with the reference pattern to measure the optical
distortion.

The distortion map is written as a multi-extension FITS file containing
16 extensions:

\begin{itemize}

\item Extensions 1--4 contain the distortion map, bad-pixel map, RMS map,
and weight map for the left field of view in the detector
$x$ direction.

\item Extensions 5--8 contain the corresponding products for the left field
of view in the detector $y$ direction.

\item Extensions 9--12 contain the distortion map, bad-pixel map, RMS map,
and weight map for the right field of view in the detector
$x$ direction.

\item Extensions 13--16 contain the corresponding products for the right field
of view in the detector $y$ direction.

\end{itemize}

Optional quality-control products include images of the detected point
patterns, distortion-corrected detector and field-of-view images, and residual
distortion maps. For a successful solution, the residual distortion is
typically below 0.1 pixels.

The estimated optical-axis position and the polynomial coefficients describing
the distortion solution are recorded in the FITS header for quality monitoring
and verification.

\paragraph{Customization}

Recipe parameters:

\begin{itemize}

\item \texttt{ird.distortion\_map.coll\_alg:} 
Collapse algorithm
(\texttt{0} = Mean,
 \texttt{1} = Median,
 \texttt{2} = Clean Mean; default).

\item \texttt{ird.distortion\_map.threshold:} 
Detection threshold for point sources (default: \texttt{3.0}).

\end{itemize}


%%%%%%%%%%%%%%%%%%%%%%%%%%%%%%%%%%%%%%%%%%%%%%%%%%%%%%%%%%%%%%
\subsection{Subworkflow: Standard Imaging}

This subworkflow processes photometric and astrometric standard-star
observations using the following tasks:

\begin{itemize}
    \item \textbf{irdis\_standard\_flux\_imaging}
    \item \textbf{irdis\_standard\_astrometry\_imaging}
\end{itemize}

These tasks execute the pipeline recipes
\texttt{sph\_ird\_science\_imaging} (classical imaging, CI) 

and
\texttt{sph\_ird\_science\_dbi} (dual-band imaging, DBI), respectively, described hereafter.

%%%%%%%%%%%%%%%%%%%%%%%%%%%%%%%%%%%%%%%%%%%%%%%%%%%%%%%%%%%%%%
\subsubsection{Standard Flux Imaging}

Recipe: \texttt{sph\_ird\_science\_imaging}

The task \textbf{irdis\_standard\_flux\_imaging} uses the recipe to reduce photometric standard-star observations for zero-point calibration.

Raw frames are dark-subtracted, optionally flat-fielded, and a combined bad-pixel map is generated. The left and right detector regions are
extracted using the IRDIS instrument model. Individual exposures are then
combined using one of the available collapse algorithms.

The output product is a multi-extension FITS file containing eight extensions:
the image, bad-pixel map, weight map, and RMS map for both the left and
right fields of view.

\paragraph{Customization}

Recipe parameters:

\begin{itemize}

\item \texttt{ird.science\_imaging.coll\_alg:} 
Collapse algorithm
(\texttt{0} = Mean,
 \texttt{1} = Median,
 \texttt{2} = Clean Mean; default).

\end{itemize}

%%%%%%%%%%%%%%%%%%%%%%%%%%%%%%%%%%%%%%%%%%%%%%%%%%%%%%%%%%%%%%
\subsubsection{Standard Astrometry Imaging}

Recipe: \texttt{sph\_ird\_science\_dbi}

The task \textbf{irdis\_standard\_astrometry\_imaging} uses the recipe to reduce astrometric standard observations.

The processing is identical to that of the photometric standard recipe. Raw
frames are dark-subtracted, optionally flat-fielded, combined with an updated
bad-pixel map, and split into left and right detector regions using the IRDIS
instrument model. The individual exposures are combined using the selected
collapse algorithm.

The final product is a multi-extension FITS file containing the image,
bad-pixel map, weight map, and RMS map for both detector fields of view.

\paragraph{Customization}

Recipe parameters:

\begin{itemize}

\item \texttt{ird.science\_dbi.coll\_alg:}
Collapse algorithm
(\texttt{0} = Mean,
 \texttt{1} = Median,
 \texttt{2} = Clean Mean; default).

\end{itemize}

%%%%%%%%%%%%%%%%%%%%%%%%%%%%%%%%%%%%%%%%%%%%%%%%%%%%%%%%%%%%%%
\subsection{Subworkflow: On-Sky Science Calibrations}

Recipes: 
\texttt{sph\_ird\_science\_dpi},
\texttt{sph\_ird\_science\_dbi},
\texttt{sph\_ird\_star\_center}

This subworkflow processes on-sky calibration observations acquired in IRDIS
imaging (classical, dual band) and polarimetry modes.

%%%%%%%%%%%%%%%%%%%%%%%%%%%%%%%%%%%%%%%%%%%%%%%%%%%%%%%%%%%%%%
\subsubsection{Imaging Flux Standard}

The task \textbf{irdis\_science\_flux\_imaging} executes the recipe \texttt{sph\_ird\_science\_dbi} to reduce flux-calibration frames acquired for coronagraphic observations, in which the target star is intentionally offset outside the coronagraphic mask.

These products are intended for quality control only and are not used during the science reduction, since no absolute flux calibration is applied to the final science products.

%%%%%%%%%%%%%%%%%%%%%%%%%%%%%%%%%%%%%%%%%%%%%%%%%%%%%%%%%%%%%%
\subsubsection{Polarimetry Flux Standard}

The task \textbf{irdis\_science\_flux\_polarimetry} executes the recipe \texttt{sph\_ird\_science\_dpi} to reduce flux-calibration observations acquired for coronagraphic polarimetric observations.

As for imaging mode, these products are not propagated into the science reduction and serve only as quality-control products.

%%%%%%%%%%%%%%%%%%%%%%%%%%%%%%%%%%%%%%%%%%%%%%%%%%%%%%%%%%%%%%
\subsubsection{Coronagraph Centre}

Recipe: \texttt{sph\_ird\_star\_center}

The task \textbf{irdis\_coronagraph\_center} determines the stellar position
behind the coronagraph by analysing waffle calibration images.

The illuminated left and right detector regions are processed independently
using a source-detection algorithm with a user-defined detection threshold.
For each exposure, the output table records:

\begin{itemize}
    \item the exposure start time,
    \item the measured stellar centre coordinates,
    \item the IRDIS DMS position converted from microns to detector pixels.
\end{itemize}

The output consists of a binary FITS table containing the measured centre
positions for all processed frames.

\paragraph{Customization}

Recipe parameters:

\begin{itemize}

\item \texttt{ird.star\_center.coll\_alg:} 
Collapse algorithm
(\texttt{0} = Mean,
 \texttt{1} = Median,
 \texttt{2} = Clean Mean; default).

\item \texttt{ird.star\_center.sigma:} 
Sigma threshold used for source detection
(default: \texttt{10.0}).

\end{itemize}

%%%%%%%%%%%%%%%%%%%%%%%%%%%%%%%%%%%%%%%%%%%%%%%%%%%%%%%%%%%%%%
\subsection{Subworkflow: Wavelength Calibration}

Recipe: \texttt{sph\_ird\_wave\_calib}

In long-slit spectroscopy (IRDIS-LSS) mode, the task
\textbf{irdis\_wavecal\_spectroscopy} derives the wavelength calibration using
the recipe \texttt{sph\_ird\_wave\_calib}.

The input calibration frames are combined, dark-subtracted, flat-fielded, and
cleaned of bad pixels. The resulting image is sliced along the dispersion
direction, and calibration emission lines are searched for within a window of
$\pm$\texttt{ird.wave\_calib.line\_tolerance} around their expected detector
positions.

The measured line positions are matched to their reference wavelengths and
fitted with a polynomial. This polynomial defines the wavelength solution and
provides the interpolation between calibration lines. The final solution is
written to the Pixel Description Table (PDT).

The output product is a FITS file containing six image extensions
corresponding to the PDT columns:

\begin{itemize}
    \item wavelength,
    \item spectrum identifier,
    \item slit identifier,
    \item wavelength width (or uncertainty),
    \item second derivative,
    \item illumination fraction.
\end{itemize}

Additional extensions contain the combined calibration image, a bad-pixel map,
and an RMS map.

The wavelength-calibration task produces \texttt{IRD\_WAVECALIB} products.
These products are associated with spectroscopy science observations for
calibration-selection purposes only and are filtered out from the recipe input
of \texttt{sph\_ird\_science\_spectroscopy}.

\paragraph{Customization}

Recipe parameters:

\begin{itemize}

\item \texttt{ird.wave\_calib.coll\_alg:}
Collapse algorithm
(\texttt{0} = Mean,
 \texttt{1} = Median,
 \texttt{2} = Clean Mean; default).

\end{itemize}



%%%%%%%%%%%%%%%%%%%%%%%%%%%%%%%%%%%%%%%%%%%%%%%%%%%%%%%%%%%%%%
\subsection{Science Imaging}

Recipe: \texttt{sph\_ird\_science\_dbi}

The task \textbf{irdis\_science\_imaging} reduces IRDIS imaging science observations, including classical, dual-band, coronagraphic, and SAM modes, using the recipe \texttt{sph\_ird\_science\_dbi}.
The recipe supports observations obtained with
or without dithering and optionally performs Angular Differential Imaging (ADI)
and Spectral Differential Imaging (SDI) processing.

Raw frames are dark-subtracted, optionally flat-fielded, and
combined with the static bad-pixel information derived from the calibration
products. The left and right detector regions are extracted using the IRDIS
instrument model before the individual exposures are combined using the selected
collapse algorithm.

A dark or background calibration frame is required. If several suitable
products are available, they are selected according to the following priority:
\texttt{SKY\_BG}
$\rightarrow$
\texttt{INS\_BG}
$\rightarrow$
\texttt{MASTER\_DARK}.

The output is a multi-extension FITS file containing eight extensions, namely, for each detector image (right and left): the science image, the bad-pixel map, the weight map (ncomb), and the RMS map.

When the required calibration tables are available, the recipe can also compute quality-control parameters such as the Strehl ratio.

\paragraph{Customization}

Recipe parameters:

\begin{itemize}

\item \texttt{ird.science\_dbi.coll\_alg:}
Collapse algorithm
(\texttt{0} = Mean,
 \texttt{1} = Median,
 \texttt{2} = Clean Mean; default).

\item \texttt{ird.science\_dbi.use\_adi:}
Enables Angular Differential Imaging processing
(\texttt{0} = disabled; default,
 \texttt{1} = enabled).

\item \texttt{ird.science\_dbi.use\_sdi:}
Enables Spectral Differential Imaging processing
(\texttt{0} = disabled; default,
 \texttt{1} = enabled).

\end{itemize}

%%%%%%%%%%%%%%%%%%%%%%%%%%%%%%%%%%%%%%%%%%%%%%%%%%%%%%%%%%%%%%
\subsection{Science Polarimetry}

Recipe: \texttt{sph\_ird\_science\_dpi}

The task \textbf{irdis\_science\_polarimetry} reduces science observations
acquired in IRDIS Dual-Polarization Imaging (DPI) mode using the pipeline
recipe \texttt{sph\_ird\_science\_dpi}.

The reduction procedure is essentially identical to that of the DBI science
recipe. Raw frames are dark-subtracted, optionally flat-fielded, corrected for
bad pixels, and combined using the selected collapse algorithm. The principal
difference lies in the format of the final products, which are provided as
polarization and intensity images rather than left and right detector images.

The standard output is a multi-extension FITS file containing eight extensions:

\begin{itemize}
    \item Intensity image:
    \begin{itemize}
        \item science image,
        \item bad-pixel map,
        \item weight map (ncomb),
        \item RMS map.
    \end{itemize}
    
    \item Polarization image:
    \begin{itemize}
        \item science image,
        \item bad-pixel map,
        \item weight map (ncomb),
        \item RMS map;
    \end{itemize}
\end{itemize}

When ADI processing is enabled, only the polarization products are written,
resulting in a FITS file with four extensions.

\paragraph{Customization}

Recipe parameters:

\begin{itemize}

\item \texttt{ird.science\_dpi.coll\_alg:}
Collapse algorithm
(\texttt{0} = Mean,
 \texttt{1} = Median,
 \texttt{2} = Clean Mean; default).

\item \texttt{ird.science\_dpi.use\_adi:}
Enables Angular Differential Imaging processing
(\texttt{0} = disabled; default,
 \texttt{1} = enabled).

\item \texttt{ird.science\_dpi.use\_sdi:}
Enables Spectral Differential Imaging processing
(\texttt{0} = disabled; default,
 \texttt{1} = enabled).

\end{itemize}

%%%%%%%%%%%%%%%%%%%%%%%%%%%%%%%%%%%%%%%%%%%%%%%%%%%%%%%%%%%%%%
\subsection{Science Spectroscopy}

Recipe: \texttt{sph\_ird\_science\_spectroscopy}

The task \textbf{irdis\_science\_spectroscopy} reduces science observations obtained in IRDIS Long-Slit Spectroscopy (LSS) mode using the pipeline recipe
\texttt{sph\_ird\_science\_spectroscopy}.

The reduction consists of dark subtraction, optional flat-field correction, bad
pixel handling, and combination of the individual exposures using the selected
collapse algorithm. EDPS subtracts an associated background calibration (sky background preferred, otherwise instrumental background, otherwise master dark).

The output product is a multi-extension FITS file containing the reduced
two-dimensional spectrum for one detector field of view (left or right,
depending on the selected slit). The FITS file contains four extensions: the science image, the bad-pixel map, the weight map (ncomb), and the RMS map.

For observations in which both left and right spectra are recorded, separate
products are generated for each detector field of view.

\paragraph{Customization}

Recipe parameters:

\begin{itemize}

\item \texttt{ird.science\_spectroscopy.coll\_alg:}
Collapse algorithm
(\texttt{0} = Mean,
 \texttt{1} = Median,
 \texttt{2} = Clean Mean; default).

\end{itemize}
